\documentclass[11pt,a4paper]{article}

\usepackage[margin=2.5cm]{geometry}
\usepackage{amsmath,amssymb,amsthm}
\usepackage{booktabs}
\usepackage{enumitem}
\usepackage{hyperref}
\usepackage{microtype}

\newtheorem{theorem}{Theorem}
\newtheorem{proposition}[theorem]{Proposition}
\newtheorem{lemma}[theorem]{Lemma}
\newtheorem{corollary}[theorem]{Corollary}
\theoremstyle{definition}
\newtheorem{definition}[theorem]{Definition}
\newtheorem{example}[theorem]{Example}
\theoremstyle{remark}
\newtheorem{remark}[theorem]{Remark}

\newcommand{\Z}{\mathbb{Z}}
\newcommand{\Q}{\mathbb{Q}}
\newcommand{\N}{\mathbb{N}}
\newcommand{\Gal}{\operatorname{Gal}}
\newcommand{\Fix}{\operatorname{Fix}}
\newcommand{\Frob}{\operatorname{Frob}}
\newcommand{\li}{\operatorname{li}}

\title{Galois Theory and the Shield Bound:\\Can Any Polynomial Beat Euler's $\boldsymbol{p^2+p+41}$?}
\author{Robin Langer}
\date{August 2025}

\begin{document}
\maketitle

\begin{abstract}
Euler's polynomial $p^2+p+41$ produces primes at remarkable density: 44.3\% of its values at prime arguments $p < 10{,}000$ are prime.
We use the Bateman--Horn conjecture and the Chebotarev density theorem to determine what Galois-theoretic conditions a polynomial must satisfy to beat Euler's asymptotic prime production rate.\looseness=-1
We show that no polynomial of degree $\geq 3$ can compete: the ``degree tax'' in Bateman--Horn requires cubics to block all primes up to $\sim\!251$ (quartics up to $\sim\!1759$), which Chebotarev renders infeasible.
Among quadratics, we find that Euler's discriminant $\Delta = -163$ is \emph{not} optimal: the composite discriminant $\Delta = -85507 = -37 \cdot 2311$ achieves a higher Bateman--Horn constant by exploiting ramification to extend the prime-blocking shield.
This reveals that the classical Heegner number story is incomplete --- composite discriminants with strategically placed ramification can outperform the class-number-one miracle.
\end{abstract}

\tableofcontents

\section{Introduction}\label{sec:intro}

The polynomial $f(x) = x^2 + x + 41$ is one of the most celebrated objects in elementary number theory.
Discovered by Euler, it takes prime values for $x = 0, 1, \ldots, 39$, a streak explained by the fact that its discriminant $\Delta = -163$ is the largest Heegner number, corresponding to the imaginary quadratic field $\Q(\sqrt{-163})$ of class number one.

When we evaluate $f$ at prime arguments $p < 10{,}000$, it produces primes 44.3\% of the time (545 out of 1229 primes).
The Bateman--Horn conjecture predicts that this dominance persists asymptotically, with a constant $C_f/\deg = 3.34$ governing the prime production rate.

The natural question is: \emph{can any polynomial beat Euler?}
We bring Galois theory --- specifically the Chebotarev density theorem and the fixed-point statistics of Galois groups --- to bear on this question.

\subsection{Summary of results}

\begin{enumerate}[label=(\roman*)]
\item For each degree $d$ and transitive Galois group $G \leq S_d$, we compute the minimum \emph{shield bound} $B$ --- the largest prime up to which \emph{all} primes must be blocked (have no roots mod $p$) --- required to beat Euler. For $d=3$ this is $B \approx 263$; for $d=4$ it is $B \approx 1759$; for $d \geq 5$ it exceeds $10{,}000$. (Section~\ref{sec:shield}.)

\item The probability that a random polynomial with Galois group $G$ achieves shield bound $B$ is roughly $\delta(G)^{\pi(B)}$, where $\delta(G)$ is the derangement fraction. For $G = \Z/3\Z$ this is $(2/3)^{56} \approx 10^{-10}$; for all degree-4 and degree-5 groups it is below $10^{-35}$. (Section~\ref{sec:shield}.)

\item An exhaustive search over negative discriminants $\Delta$ down to $-100{,}000$ reveals three quadratics that beat Euler on the Bateman--Horn constant $C_f/2$. The winner is $\Delta = -85507 = -37 \cdot 2311$ with $C_f/2 = 3.48$. (Section~\ref{sec:quadratic}.)

\item The mechanism is \emph{ramification shielding}: the prime $37$ divides $\Delta$, making it ramified (contributing a neutral factor of~1 to the Bateman--Horn product), while all primes from 2 to 59 are quadratic non-residues of the cofactor. (Section~\ref{sec:ramification}.)

\item The best cyclic cubic found (conductor $q = 4{,}851{,}139$, blocking all primes up to 137) achieves $C_f/3 = 2.88$, still 14\% behind Euler. (Section~\ref{sec:cubic}.)
\end{enumerate}


\section{The Bateman--Horn framework}\label{sec:bh}

Let $f \in \Z[x]$ be an irreducible polynomial of degree $d \geq 2$ with positive leading coefficient.
The \emph{Bateman--Horn conjecture}~\cite{BH} predicts
\[
\pi_f(N) \;=\; \#\{n \leq N : f(n) \text{ is prime}\} \;\sim\; \frac{C_f}{d} \cdot \li(N),
\]
where $\li(N) = \int_2^N dt/\ln t$ is the logarithmic integral and
\begin{equation}\label{eq:BH}
C_f \;=\; \prod_{p \text{ prime}} \frac{1 - \omega_f(p)/p}{1 - 1/p},
\end{equation}
with $\omega_f(p) = \#\{x \in \Z/p\Z : f(x) \equiv 0 \pmod{p}\}$.

The key observation is that $C_f/d$ governs the asymptotic density.
To beat Euler, a polynomial of degree $d$ needs $C_f/d > 3.34$.

\begin{definition}[Shield]
A prime $p$ is \emph{blocked} by $f$ if $\omega_f(p) = 0$ (i.e., $f$ has no roots mod $p$).
The \emph{shield} of $f$ is the set of blocked primes.
The \emph{shield product} is $\prod_{p \text{ blocked}} \frac{p}{p-1}$.
\end{definition}

When $\omega_f(p) = 0$, the factor at $p$ in~\eqref{eq:BH} is $\frac{p}{p-1} > 1$, boosting $C_f$.
When $\omega_f(p) = k > 0$, the factor is $\frac{p-k}{p-1} < 1$, penalising $C_f$.

\begin{proposition}[Average neutrality]\label{prop:neutral}
For any transitive subgroup $G \leq S_d$, the expected value of the Bateman--Horn factor at a prime $p$ (averaged over Frobenius elements by Chebotarev) equals~$1$.
\end{proposition}

\begin{proof}
The expected number of fixed points of a uniformly random element of a transitive group is exactly~$1$ (a standard result in permutation group theory).
Therefore
\[
\mathbb{E}\!\left[\frac{1 - \omega/p}{1 - 1/p}\right]
= \frac{1}{1-1/p}\left(1 - \frac{\mathbb{E}[\omega]}{p}\right)
= \frac{1}{1-1/p}\cdot\frac{p-1}{p}
= 1. \qedhere
\]
\end{proof}

This is fundamental: $C_f$ is not determined by the ``average'' blocking rate (which is always neutral) but by \emph{which specific small primes} are blocked, since $p/(p-1)$ is largest for small $p$.


\section{Shield bounds by Galois group}\label{sec:shield}

For a polynomial $f$ of degree $d$ with Galois group $G$ acting transitively on the $d$ roots, the Chebotarev density theorem determines the distribution of $\omega_f(p)$ over primes:
\[
\omega_f(p) = |\Fix(\Frob_p)|,
\]
where $\Frob_p \in G$ is the Frobenius element at $p$ (well-defined up to conjugacy for unramified $p$).

The \emph{derangement fraction} $\delta(G) = |\{g \in G : \Fix(g) = \varnothing\}|/|G|$ gives the asymptotic proportion of blocked primes.

\begin{table}[ht]
\centering
\caption{Shield bounds required to beat Euler ($C_f/d > 3.34$).}
\label{tab:shield}
\medskip
\begin{tabular}{llcccc}
\toprule
Group $G$ & Degree & $\delta(G)$ & Need $C_f >$ & Shield $B$ & Feasibility \\
\midrule
$\Z/2$ & 2 & 50.0\% & 6.68 & 37 & Euler achieves this \\
$\Z/3$ & 3 & 66.7\% & 10.02 & 263 & $\sim\!10^{10}$ candidates \\
$S_3$  & 3 & 33.3\% & 10.02 & 263 & infeasible ($\sim\!10^{27}$) \\
$\Z/4$ & 4 & 75.0\% & 13.36 & 1759 & infeasible ($\sim\!10^{35}$) \\
$V_4$  & 4 & 75.0\% & 13.36 & 1759 & infeasible ($\sim\!10^{35}$) \\
$D_4$  & 4 & 62.5\% & 13.36 & 1759 & infeasible ($\sim\!10^{56}$) \\
$\Z/5$ & 5 & 80.0\% & 16.70 & 11789 & infeasible ($\sim\!10^{137}$) \\
$F_{20}$ & 5 & 75.0\% & 16.70 & 11789 & infeasible ($\sim\!10^{177}$) \\
\bottomrule
\end{tabular}
\end{table}

The shield bound $B$ is the smallest prime such that $\prod_{p \leq B} \frac{p}{p-1} > 3.34 \cdot d$.
By Mertens' theorem, this product grows as $e^\gamma \ln B \approx 1.781 \ln B$.

The ``feasibility'' column estimates the number of candidates that must be searched: roughly $\delta(G)^{-\pi(B)}$, the reciprocal of the probability that all primes up to $B$ are simultaneously blocked.

\begin{corollary}[The degree tax]
For the same shield (blocking all primes up to $B$), a degree-$d$ polynomial has $C_f/d \approx \frac{1}{d}\prod_{p \leq B}\frac{p}{p-1}$.
The factor $1/d$ is an inescapable penalty that grows linearly with degree.
\end{corollary}


\section{Euler and the Heegner miracle}\label{sec:euler}

For $f(x) = x^2 + x + 41$, the discriminant is $\Delta = 1 - 4 \cdot 41 = -163$.

A prime $p$ is blocked if and only if $\left(\frac{\Delta}{p}\right) = -1$, i.e., $\Delta$ is a quadratic non-residue mod $p$.
By quadratic reciprocity, for odd $p \neq 163$:
\[
\left(\frac{-163}{p}\right) = \left(\frac{p}{163}\right),
\]
so $p$ is blocked iff $p$ is a QNR mod $163$.

The field $\Q(\sqrt{-163})$ has class number $h(-163) = 1$.
By the Rabinowitsch criterion, this is equivalent to $x^2 + x + 41$ being prime for all $x = 0, 1, \ldots, 39$.
All primes $p < 41$ are QNRs mod $163$, giving a shield of 12 consecutive primes.

The shield product is
\[
\prod_{\substack{p \leq 37}} \frac{p}{p-1} = \frac{2}{1}\cdot\frac{3}{2}\cdot\frac{5}{4}\cdots\frac{37}{36} \approx 6.72.
\]
Beyond $p = 41$, about half of primes remain blocked (by Chebotarev), contributing additional small corrections.
The final Bateman--Horn constant is $C_f = 6.68$, giving $C_f/2 = 3.34$.

This connects to deep mathematics:
\begin{itemize}
\item Ramanujan's constant: $e^{\pi\sqrt{163}} \approx 640320^3 + 744$ (almost an integer),
\item The $j$-invariant: $j\!\left(\frac{1+\sqrt{-163}}{2}\right) = -640320^3$,
\item Stark--Heegner: there are exactly nine imaginary quadratic fields with class number~$1$.
\end{itemize}


\section{Beating Euler: composite discriminants}\label{sec:quadratic}

We searched all negative discriminants $\Delta \equiv 0, 1 \pmod{4}$ down to $-100{,}000$, computing the Bateman--Horn constant for each.

\begin{table}[ht]
\centering
\caption{Top quadratics by Bateman--Horn constant $C_f/2$.}
\label{tab:top-quad}
\medskip
\begin{tabular}{rrrrcl}
\toprule
Rank & $\Delta$ & $C_f$ & $C_f/2$ & 1st split & Polynomial \\
\midrule
1 & $-85507$ & 6.964 & 3.482 & 61 & $x^2+x+21377$ \\
2 & $-90787$ & 6.778 & 3.389 & 37 & $x^2+x+22697$ \\
3 & $-77683$ & 6.705 & 3.353 & 47 & $x^2+x+19421$ \\
4 & $-163$   & 6.682 & 3.341 & 41 & $x^2+x+41$ \\
5 & $-86227$ & 6.391 & 3.196 & 41 & $x^2+x+21557$ \\
\bottomrule
\end{tabular}
\end{table}

The winner, $\Delta = -85507$, beats Euler by 4.2\% on $C_f/2$.


\subsection{The ramification mechanism}\label{sec:ramification}

The discriminant $-85507$ factors as $-37 \cdot 2311$, with both $37$ and $2311$ prime.

\begin{itemize}
\item $p = 37$: since $37 \mid \Delta$, the prime is \emph{ramified} in $\Q(\sqrt{-85507})$, giving $\omega_f(37) = 1$ and a Bateman--Horn factor of $\frac{1-1/37}{1-1/37} = 1$ --- completely neutral.

\item For all primes $p < 61$ with $p \neq 37$: the Kronecker symbol $\left(\frac{-85507}{p}\right) = -1$, so $\omega_f(p) = 0$ (blocked).
\end{itemize}

This creates a ``shield with a neutral hole'': the 16 primes $\{2, 3, 5, 7, 11, 13, 17, 19, 23, 29, 31, 41, 43, 47, 53, 59\}$ are blocked, while $37$ sits neutrally in the middle.
The shield product is $\approx 8.13$, compared to Euler's $\approx 7.89$.

\begin{remark}
Euler's polynomial benefits from class number $h(-163) = 1$, which guarantees all primes up to $\sqrt{163/3} \approx 7.4$ are blocked, with the full shield extending to 37 by the specific arithmetic of $163$.
The discriminant $-85507$ has class number $> 1$, but its shield extends further by exploiting ramification to ``skip over'' the split prime at $37$.
\end{remark}


\subsection{Crossover analysis}

Although $\Delta = -85507$ has a higher Bateman--Horn constant, the polynomial $x^2+x+21377$ produces larger values (the constant term is 520 times bigger), making individual values less likely to be prime at moderate arguments.

\begin{table}[ht]
\centering
\caption{Prime counts at prime arguments: Euler vs.\ $\Delta = -85507$.}
\label{tab:crossover}
\medskip
\begin{tabular}{rcccc}
\toprule
$N$ & Euler primes & $-85507$ primes & Euler density & $-85507$ density \\
\midrule
1\,000 & 109 & 110 & 0.649 & 0.655 \\
2\,000 & 168 & 176 & 0.555 & 0.581 \\
5\,000 & 329 & 320 & 0.492 & 0.478 \\
10\,000 & 545 & 531 & 0.443 & 0.432 \\
20\,000 & 917 & 907 & 0.405 & 0.401 \\
50\,000 & 1885 & 1901 & 0.367 & 0.370 \\
\bottomrule
\end{tabular}
\end{table}

The two polynomials trade the lead repeatedly.
By $N = 50{,}000$, the Bateman--Horn prediction begins to dominate: $\Delta = -85507$ pulls ahead with 1901 vs.\ 1885 primes.
The asymptotic advantage of $\sim\!4\%$ should become decisive for large $N$.


\section{The cyclic cubic attack}\label{sec:cubic}

The most promising non-quadratic avenue is cyclic cubics ($G = \Z/3$), which block $2/3$ of primes --- the highest derangement fraction among degree-3 groups.

For a cyclic cubic with prime conductor $q$ (where $3 \mid q-1$), the splitting of a prime $p$ is controlled by cubic reciprocity:
\[
\omega_f(p) = \begin{cases} 3 & \text{if } p^{(q-1)/3} \equiv 1 \pmod{q} \text{ (cubic residue)}, \\ 0 & \text{otherwise}. \end{cases}
\]

We searched all primes $q \equiv 1 \pmod{3}$ up to $10^7$, looking for conductors where all small primes are cubic non-residues.

\begin{table}[ht]
\centering
\caption{Best cyclic cubic conductors.}
\label{tab:cubic}
\medskip
\begin{tabular}{rccccc}
\toprule
Conductor $q$ & 1st cubic residue & \# blocked & Shield & $C_f/3$ & Density \\
\midrule
$4{,}851{,}139$ & 137 & 32 & 8.850 & 2.878 & 0.339 \\
$7{,}778{,}509$ & 131 & 31 & 8.782 & 2.862 & 0.359 \\
$6{,}843{,}919$ & 127 & 30 & 8.713 & 2.934 & 0.356 \\
$2{,}206{,}591$ & 113 & 29 & 8.636 & 2.964 & 0.371 \\
\bottomrule
\end{tabular}
\end{table}

The best cubic blocks all primes up to 137 --- an impressive shield --- but $C_f/3 \approx 2.9$ remains 14\% below Euler's $3.34$.

For each conductor, we constructed explicit cubic polynomials via Gauss periods.
For conductor $q$, let $g$ be a primitive root mod $q$ and partition $(\Z/q\Z)^\times$ into three cosets of the subgroup of cubes.
The Gauss period $\eta_0 = \sum_{k \in C_0} \cos(2\pi k/q)$ is a root of a monic cubic with integer coefficients.

For example, conductor $q = 2{,}206{,}591$ yields $x^3 + x^2 - 735{,}530x + 227{,}687{,}501$, a cyclic cubic blocking all 25 primes up to~97.


\section{Why quadratics win}\label{sec:why}

We can now state the structural result:

\begin{theorem}[Informal]
Under the Bateman--Horn conjecture, no polynomial of degree $d \geq 3$ can match the asymptotic prime production rate of the optimal quadratic.
Among quadratics, the optimal discriminant $($by Bateman--Horn constant$)$ is not a Heegner number but a composite discriminant exploiting ramification.
\end{theorem}

The argument rests on three pillars:

\paragraph{1. The degree tax.}
The Bateman--Horn prediction scales as $C_f/d$.
For the same shield product $S$, a degree-$d$ polynomial achieves $C_f/d \approx S/d$.
To match Euler's $S/2 \approx 3.34$, degree $d$ needs $S > 3.34 \cdot d$, which by Mertens' theorem requires blocking all primes up to $\exp(3.34d / e^\gamma) \approx \exp(1.88d)$.
This grows exponentially with $d$.

\paragraph{2. Chebotarev makes large shields improbable.}
By Proposition~\ref{prop:neutral}, only a fraction $\delta(G)$ of primes are blocked.
The probability that \emph{all} primes up to $B$ are blocked is roughly $\delta(G)^{\pi(B)}$, which decreases super-exponentially with $B$.
Even the best group ($\Z/5$ with $\delta = 4/5$) cannot achieve the shield needed at degree~5.

\paragraph{3. Ramification gives quadratics an extra trick.}
Composite discriminants allow specific primes to be ramified (neutral) rather than split (penalised), effectively ``skipping over'' bad primes in the shield.
This trick is unique to the interplay of the discriminant's factorisation with quadratic reciprocity, and has no direct analogue for cubic or higher reciprocity that would overcome the degree tax.


\section{Open questions}\label{sec:open}

\begin{enumerate}
\item \textbf{Optimal composite discriminants.}
Is there a theoretical framework, analogous to Heegner's theorem for class number one, that predicts which composite discriminants maximise $C_f$?
The interplay of ramification and quadratic reciprocity at composite moduli is rich and, to our knowledge, not well characterised for this optimisation problem.

\item \textbf{Deeper cubic search.}
Our search covered conductors up to $10^7$, finding shields up to $B = 137$.
A search to $10^{10}$ might find shields up to $\sim\!200$, but $C_f/3$ would need to exceed $3.48$ (the new quadratic champion), requiring a shield product above $10.4$ --- blocking all primes up to $\sim\!290$.

\item \textbf{Function field analogue.}
Over $\mathbb{F}_q[t]$, the Bunyakovsky conjecture is a theorem (via Weil's proof of the Riemann hypothesis for curves).
Does the proof structure illuminate which polynomials over $\Z$ should produce the most primes?

\item \textbf{Non-monic quadratics.}
We tested $f(x) = ax^2 + 1$ for various primorial values of $a$.
Even $a = 30030 = 2 \cdot 3 \cdot 5 \cdot 7 \cdot 11 \cdot 13$ achieves only $C_f/2 = 2.26$.
The ``free'' blocking of primes dividing $a$ cannot compensate for the poor residual discriminant $-4a$.
A more systematic search over general $ax^2 + bx + c$ might be worthwhile.
\end{enumerate}

\begin{thebibliography}{9}
\bibitem{BH}
P.~T.~Bateman and R.~A.~Horn,
\emph{A heuristic asymptotic formula concerning the distribution of prime numbers},
Math.\ Comp.\ \textbf{16} (1962), 363--367.

\bibitem{Heegner}
K.~Heegner,
\emph{Diophantische Analysis und Modulfunktionen},
Math.\ Z.\ \textbf{56} (1952), 227--253.

\bibitem{Rabinowitsch}
G.~Rabinowitsch,
\emph{Eindeutigkeit der Zerlegung in Primzahlfaktoren in quadratischen Zahlk\"orpern},
J.\ Reine Angew.\ Math.\ \textbf{142} (1913), 153--164.

\bibitem{Stark}
H.~M.~Stark,
\emph{A complete determination of the complex quadratic fields of class-number one},
Michigan Math.\ J.\ \textbf{14} (1967), 1--27.
\end{thebibliography}

\end{document}
